\documentclass[11pt,letterpaper]{article}
\usepackage[margin=0.8in]{geometry}
\usepackage[T1]{fontenc}
\usepackage[utf8]{inputenc}
\usepackage{lmodern}
\usepackage{amsmath,amssymb,booktabs,array,graphicx}
\usepackage[hidelinks]{hyperref}
\usepackage{enumitem}
\DeclareUnicodeCharacter{2212}{\ensuremath{-}}
\DeclareUnicodeCharacter{FF1A}{:}
\setcounter{secnumdepth}{0}
\setlength{\parindent}{0pt}
\setlength{\parskip}{6pt}
\setlist{nosep,leftmargin=*}
\providecommand{\tightlist}{\setlength{\itemsep}{0pt}\setlength{\parskip}{0pt}}
\emergencystretch=2em
\begin{document}
\begin{center}
{\Large\bfseries Problem Set 5: Questions (c)--(e)}
\end{center}

\hypertarget{question-c-report}{%
\subsection{Question (c) Report}\label{question-c-report}}

\textbf{How the PS3 code was extended.} I used an LLM to produce most of the PS5 code extension, starting from my PS3 implementation in cells 1(b) and 1(f). The main changes were adding joint regressions on RMRF, SMB and HML to estimate the additional factor betas, generalizing the monthly cross-sectional regression into a function that runs all three specifications, and producing a side-by-side results table. The extension preserved the original FMB standard-error calculation, fixed betas across months, and lagged, time-varying characteristics. The final implementation passes the included verification checks and reproduces the PS3 size and BE/ME estimates to four decimal places when using the original sample and market-only beta. Equation (1) retains the PS3 market-only beta, while equations (2) and (3) use joint three-factor betas, following our stated team convention.

\textbf{Were the changes correct?} Yes. All checks in Cell 7 pass:

\begin{itemize}
\item
  Equation (2) equals a single regression of average returns on the fixed betas.
\item
  A numpy re-run of one month confirms which beta each equation uses.
\item
  All equations use the same 1,194 months with 25 portfolios each.
\item
  Every SE equals \(sd(\hat\gamma_t)/\sqrt{T}\).
\end{itemize}

\textbf{Fixed betas, time-varying characteristics.} The betas are estimated once and stay fixed across all months of the second pass. Size and BE/ME are lagged and change over time. Checks 1--3 in Cell 7 confirm this structure was preserved.

\textbf{Results} (percent per month; t-statistics test \(\gamma = 0\); 1927-01 to 2026-06, T = 1194, N = 25):

\begin{center}
\small
\resizebox{\textwidth}{!}{%
\begin{tabular}{lrrrrrrrrr}
\toprule
Coefficient & Eq (1) Est. & Eq (1) SE & Eq (1) t & Eq (2) Est. & Eq (2) SE & Eq (2) t & Eq (3) Est. & Eq (3) SE & Eq (3) t \\
\midrule
\texttt{gamma\_0} & 1.7547 & 0.4228 & 4.15 & 1.6433 & 0.3954 & 4.16 & 1.4108 & 0.3994 & 3.53 \\
\texttt{gamma\_M} & -0.4777 & 0.3563 & -1.34 & -0.9042 & 0.4095 & -2.21 & 0.2049 & 0.3911 & 0.52 \\
\texttt{gamma\_size} & -0.0674 & 0.0294 & -2.30 &  &  &  & -0.0932 & 0.0477 & -1.96 \\
\texttt{gamma\_BM} & 0.2327 & 0.0570 & 4.08 &  &  &  & 0.4808 & 0.1096 & 4.39 \\
\texttt{gamma\_SMB} &  &  &  & 0.1087 & 0.0973 & 1.12 & -0.1567 & 0.1892 & -0.83 \\
\texttt{gamma\_HML} &  &  &  & 0.3752 & 0.1061 & 3.54 & -0.4656 & 0.2270 & -2.05 \\
\bottomrule
\end{tabular}
}
\end{center}

Note: \(\beta_M\) in equation (1) is the market-only beta, and \(\beta_M\) in equations (2) and (3) is the joint FF3 loading. They are different numbers.

\textbf{Check against PS3.} On 1969-07 to 2026-06 (T = 684), with the market-only beta re-estimated on that sample, equation (1) gives \texttt{gamma\_size} = -0.0194 (t = -0.64) and \texttt{gamma\_BM} = 0.1573 (t = 2.11). These match my PS3 results exactly to 4 d.p., and \texttt{gamma\_0} and \texttt{gamma\_M} also match. The full-sample equation (1) estimates differ because the sample period is different.

\textbf{Implementation qualifications.}

\begin{itemize}
\item
  \emph{Market beta in equation (1):} the wording is ambiguous (see the sensitivity note). With \texttt{beta\_M\_3f}, \texttt{gamma\_size} changes from -0.0674 (t = -2.30) to -0.0367 (t = -1.20), while \texttt{gamma\_BM} stays at t \textgreater{} 3.
\item
  \emph{BE/ME timing:} this follows the PS3/team calendar-year rule, not the Fama--French July--June timing described in Lecture 2.
\item
  \emph{Generated regressors:} the betas are estimated on the full sample, so they contain measurement error and a look-ahead element. The FMB SEs are not adjusted for beta-estimation error.
\item
  \emph{PS3 comparison precision:} the match is confirmed to 4 d.p., which is all PS3 printed.
\end{itemize}

\hypertarget{sensitivity-check-which-market-beta-to-use-in-equation-1}{%
\subsubsection{Sensitivity check: which market beta to use in equation (1)}\label{sensitivity-check-which-market-beta-to-use-in-equation-1}}

\textbf{Wording ambiguity.} Step 1 of Question (c) says to estimate \(\beta_{iM}\), \(\beta_{iSMB}\) and \(\beta_{iHML}\) "using full-period OLS regressions on RMRF, SMB, and HML," and the same symbol \(\beta_{iM}\) appears in all three equations. This can be read as one joint FF3 regression whose market loading is used everywhere. However, the same step notes that "\(\beta_{iM}\) was already estimated in Problem Set 3," and we are asked to verify that equation (1) matches the PS3 \(\gamma_{size}\) and \(\gamma_{B/M}\). In PS3, \(\beta_{iM}\) came from a market-only regression. The assignment does not settle which reading is intended. The team convention uses the market-only beta in (1) and the joint FF3 betas in (2) and (3).

\textbf{Main specification.} Equation (1) uses \texttt{beta\_M\_capm} (market-only, the same as PS3). Under the PS3 setup (Cell 7), this reproduces the PS3 estimates and t-statistics to 4 d.p. Equations (2) and (3) use \texttt{beta\_M\_3f}, \texttt{beta\_SMB} and \texttt{beta\_HML} from one joint regression.

\textbf{What changes with \texttt{beta\_M\_3f} in equation (1).} The sample (1927-01 to 2026-06, T = 1194) and characteristics are the same; the FMB standard-error formula is unchanged:

\begin{center}
\small
\begin{tabular}{lrrrr}
\toprule
 & \texttt{beta\_M\_capm} (main) &  & \texttt{beta\_M\_3f} (sensitivity) &  \\
\midrule
 & Estimate (SE) & t & Estimate (SE) & t \\
\texttt{gamma\_0} & 1.7547 (0.4228) & 4.15 & 1.5272 (0.4066) & 3.76 \\
\texttt{gamma\_M} & -0.4777 (0.3563) & -1.34 & -0.5264 (0.4346) & -1.21 \\
\texttt{gamma\_size} & -0.0674 (0.0294) & -2.30 & -0.0367 (0.0305) & -1.20 \\
\texttt{gamma\_BM} & 0.2327 (0.0570) & 4.08 & 0.2020 (0.0566) & 3.57 \\
\bottomrule
\end{tabular}
\end{center}

Average monthly \(R^2\) moves from 0.4898 to 0.4852, and the two market betas have a cross-sectional correlation of 0.65. The main difference is \texttt{gamma\_size}: its \textbar t\textbar{} is above 1.96 with \texttt{beta\_M\_capm} but not with \texttt{beta\_M\_3f}. \texttt{gamma\_BM} has t above 3 in both versions.

\hypertarget{question-d}{%
\subsection{Question (d)}\label{question-d}}

\hypertarget{step-1}{%
\subsubsection{Step 1}\label{step-1}}

Equation (3) has the highest average monthly cross-sectional \((R^2\)) (0.573, versus 0.490 for equation (1) and 0.485 for equation (2)), although this measures in-sample fit and partly reflects the additional regressors, so it does not by itself distinguish between the characteristics and risk explanations. In the combined specification, BE/ME stays strongly significant, and its coefficient roughly doubles relative to equation (1) (0.4808, \(t = 4.39\), versus 0.2327, \(t = 4.08\)). Size remains negative but narrowly misses the 5\% level (\(t = -1.96\), \(p = 0.0508\)). By contrast, the HML beta premium turns from positive and significant in equation (2) (0.3752, \(t = 3.54\)) to negative in equation (3) (−0.4656, \(t = -2.05\)), and SMB is insignificant in both equations (\(t = -0.83\) in (3)). Taken together, BE/ME keeps explanatory power after controlling for factor exposures, while the betas lose their expected positive sign, which leans toward a characteristics interpretation; however, the roughly doubled standard errors in equation (3) suggest betas and characteristics are hard to separate, so I treat this conclusion as tentative.

\hypertarget{step-2-llm-discussion}{%
\subsubsection{Step 2 --- LLM discussion}\label{step-2-llm-discussion}}

\textbf{1. "What is the strongest economic argument that size and BE/ME are priced because they are characteristics rather than because they co-vary with systematic risk?"}

The strongest argument is that investors may systematically misprice firms with certain observable characteristics. For example, investors may overvalue glamorous, low-BE/ME firms and undervalue distressed or neglected high-BE/ME firms, so the later return premium reflects price correction rather than compensation for bearing systematic risk. Under this view, size and BE/ME themselves should predict returns even after controlling for factor betas.

\textbf{2. "What is the strongest economic argument that size and BE/ME premia are compensation for systematic risk --- and what would that imply about the betas on SMB and HML?"}

The risk view says small and high-BE/ME firms tend to perform poorly in certain bad economic states, so investors require higher expected returns to hold them. SMB and HML are then factor-mimicking portfolios capturing these common systematic risks. Therefore, firms with higher \(\beta_{SMB}\) or \(\beta_{HML}\) should earn higher average returns, and once those betas are controlled for, size and BE/ME should have little additional explanatory power.

\textbf{3. "In a Fama--MacBeth horse race between equations (1), (2), and (3), what empirical pattern would cleanly support the characteristics view? What pattern would support the risk view?"}

Suppose equation (1) contains characteristics, equation (2) contains factor betas, and equation (3) includes both:

\begin{itemize}
\item
  Characteristics view: size and BE/ME are significant in (1) and remain significant in (3), while the SMB/HML beta prices become weak or insignificant.
\item
  Risk view: SMB/HML betas are significant in (2) and remain significant in (3), while the size and BE/ME coefficients become weak or insignificant.
\item
  If both remain significant in (3), the evidence is mixed rather than a clean victory for either explanation.
\end{itemize}

\hypertarget{evaluation}{%
\subsubsection{Evaluation：}\label{evaluation}}

\begin{enumerate}
\def\labelenumi{\arabic{enumi}.}
\tightlist
\item
  Does the LLM correctly identify the key distinguishing prediction --- that if characteristics dominate in equation (3), it is hard to sustain a pure risk story?
\end{enumerate}

The LLM correctly identifies that characteristics remaining significant in equation (3), while beta premia weaken, would challenge a pure SMB/HML risk explanation. However, significance alone is not enough and the signs of the coefficients also matter.

\begin{enumerate}
\def\labelenumi{\arabic{enumi}.}
\setcounter{enumi}{1}
\tightlist
\item
  Does the LLM correctly note that the horse race is inherently noisy because characteristics and betas are correlated by construction (size sorts produce size beta differences)?
\end{enumerate}

The revised answer does not mention that characteristics and betas are strongly correlated because of the portfolio sorts. My notebook shows correlations of −0.974 between average log size and SMB beta and 0.948 between average log BE/ME and HML beta, making their separate effects difficult to identify.

\begin{enumerate}
\def\labelenumi{\arabic{enumi}.}
\setcounter{enumi}{2}
\tightlist
\item
  Does your empirical evidence from equations (1)--(3) more strongly support one side? Be explicit about which coefficients drive your conclusion.
\end{enumerate}

My full-sample results lean toward the characteristics view, especially for BE/ME: its coefficient increases from 0.2327 (\(t=4.08\)) in equation (1) to 0.4808 (\(t=4.39\)) in equation (3). Size remains negative but narrowly misses 5\% significance (\(p=0.0508\)), while SMB is insignificant in both equations (2) and (3). HML remains significant but reverses from positive to negative (−0.4656, \(t=-2.05\)), which conflicts with the positive premium predicted by the LLM's risk explanation. The evidence therefore favors characteristics, but the strong correlations make this conclusion tentative

\hypertarget{question-e}{%
\subsection{Question (e)}\label{question-e}}

Sample: 1963-01 to 2026-06 \textbar{} T = 762 \textbar{} N = 25 \textbar{} betas estimated on the same months

gamma\_M: spec (1) on beta\_M\_capm; specs (2), (3) on beta\_M\_3f (different numbers). \% per month.

\begin{center}
\small
\resizebox{\textwidth}{!}{%
\begin{tabular}{lrrrrrrrrrrrr}
\toprule
Coefficient & Spec (1) Estimate & Spec (1) FMB SE & Spec (1) t-stat & Spec (1) p-value & Spec (2) Estimate & Spec (2) FMB SE & Spec (2) t-stat & Spec (2) p-value & Spec (3) Estimate & Spec (3) FMB SE & Spec (3) t-stat & Spec (3) p-value \\
\midrule
\texttt{gamma\_0} & 1.2790 & 0.3223 & 3.9678 & 0.0001 & 1.1429 & 0.2583 & 4.4247 & 0.0000 & 1.3700 & 0.3828 & 3.5791 & 0.0004 \\
\texttt{gamma\_M} & -0.2247 & 0.2901 & -0.7746 & 0.4388 & -0.5261 & 0.3035 & -1.7337 & 0.0834 & -0.2827 & 0.3631 & -0.7787 & 0.4364 \\
\texttt{gamma\_size} & -0.0437 & 0.0287 & -1.5202 & 0.1289 &  &  &  &  & 0.0045 & 0.0507 & 0.0892 & 0.9289 \\
\texttt{gamma\_BM} & 0.1704 & 0.0682 & 2.5007 & 0.0126 &  &  &  &  & 0.4876 & 0.1368 & 3.5652 & 0.0004 \\
\texttt{gamma\_SMB} &  &  &  &  & 0.1228 & 0.1121 & 1.0951 & 0.2738 & 0.0527 & 0.2055 & 0.2564 & 0.7977 \\
\texttt{gamma\_HML} &  &  &  &  & 0.3364 & 0.1098 & 3.0626 & 0.0023 & -0.4949 & 0.2644 & -1.8717 & 0.0616 \\
\bottomrule
\end{tabular}
}
\end{center}

Average monthly R-squared, spec (1): 0.5254

Average monthly R-squared, spec (2): 0.5283

Average monthly R-squared, spec (3): 0.5926

Post-1963 sample: 1963-01 to 2026-06 (T = 762), with betas re-estimated on the same months (team conventions). Full sample: 1927-01 to 2026-06 (T = 1194).

\begin{center}
\small
\begin{tabular}{lrrr}
\toprule
Variable & Eq & Full: Est. (t) & Post-1963: Est. (t) \\
\midrule
\(\gamma_{size}\) & (1) & -0.0674 (-2.30) & -0.0437 (-1.52) \\
\(\gamma_{B/M}\) & (1) & 0.2327 (4.08) & 0.1704 (2.50) \\
\(\gamma_{SMB}\) & (2) & 0.1087 (1.12) & 0.1228 (1.10) \\
\(\gamma_{HML}\) & (2) & 0.3752 (3.54) & 0.3364 (3.06) \\
\(\gamma_{size}\) & (3) & -0.0932 (-1.96) & 0.0045 (0.09) \\
\(\gamma_{B/M}\) & (3) & 0.4808 (4.39) & 0.4876 (3.57) \\
\(\gamma_{SMB}\) & (3) & -0.1567 (-0.83) & 0.0527 (0.26) \\
\(\gamma_{HML}\) & (3) & -0.4656 (-2.05) & -0.4949 (-1.87) \\
\bottomrule
\end{tabular}
\end{center}

\textbf{Does the result change?} For BE/ME, the result is basically the same. In equation (3), \(\gamma_{B/M}\) is almost unchanged (0.4876, t = 3.57), while the HML beta premium stays negative and is now insignificant at 5\% (t = -1.87, p = 0.06). So the value evidence still leans toward the characteristics view. The size result changes: \(\gamma_{size}\) is no longer significant in equation (1) (t = -1.52) and is essentially zero in equation (3) (t = 0.09), and SMB is insignificant in both samples. After 1963, neither size nor its beta has a statistically significant coefficient in post-1963 combined specification, so the size dimension gives no evidence for either side.

\textbf{More or less confident?} I am more confident about BE/ME, because its coefficient in equation (3) has a similar size and stays significant in both samples. I am less confident about size, which was already borderline in part (d) and disappears after 1963. Overall the conclusion is still tentative. With fewer months, the standard errors in equation (3) are larger (e.g. \(\gamma_{HML}\) SE 0.2644 vs 0.2270), and betas and characteristics are still strongly correlated by construction of the 25 portfolios.

\textbf{Why might the sample period matter?}

\begin{itemize}
\item
  \emph{Data quality:} before 1963, book equity is not from Compustat and covers fewer firms, so BE/ME and the portfolios may be measured less precisely.
\item
  \emph{In-sample vs. out-of-sample:} the post-1963 period is close to the sample in which the size and value effects were originally documented, so the pre-1963 data works partly as an out-of-sample check.
\item
  \emph{Changing premia and risk:} the size premium seems to depend on the period. The pre-1963 sample also includes the Great Depression, when distress risk was especially important. If factor loadings change over time, full-sample fixed betas fit some periods worse than others.
\end{itemize}
\end{document}
